% s-213-p213-m001; original PDF 213, 214, 215
\Aphorism{s-213-p213-m001}{306}
Die \Sperr{Shakespeare-Bacon-Frage}: eine Bewegung, deren Ziel die Abschaffung eines unerklärlichen Naturgenius und Einsetzung des gelehrten Nebenbei-Dichters, deren Ursprung das Gehirn einer Amerikanerin ist (bei jeder Angelsächsin darf man, bis zum Beweise des Gegentheils, tiefstes kirchenfrommes Ressentiment gegen böse Buben wie Shakespeare und Byron voraussetzen), und deren „Methode“ — in der Entzifferung von Geheimschriften besteht, statt in der Beglaubigung des ungeheuren Quiproquo, das drei Jahrhunderte über die meistbesprochene Persönlichkeit der Litteratur getäuscht haben soll. Wenn man von dieser „Methode“ der Alignements, der lädirten Druckbuchstaben und des Honorabilificatitudinitatibus einige Proben gekostet hat, so „kann“ man sie schon und ist um eine zerstreuende Spielerei reicher; nichts unterhaltender als mit beliebigen Freiheiten der Interpretation aus einem beliebigen Drucktext einen beliebigen Geheimsinn herauszulesen. Wie würde man z. B. in einer freien Viertelstunde beweisen, dass Kant, nicht Goethe der Dichter des Faust ist? Man nimmt etwa den Schlusschor „Alles Vergängliche“ u. s. w. mit seinen 130 Buchstaben, fügt die Überschrift „Chorus mysticus“ hinzu und hat 144 oder 12 mal 12 Buchstaben. Will man sich die Mühe machen, in der Originalausgabe die t mit abgebrochenen Schwänzen oder die e mit offenen Schleifen herauszusuchen, um so besser: sonst kann man auch der höheren Magie entrathen und schreibt einfach jene 144 Buchstaben diagonal von links unten nach rechts oben in ein Quadrat von 12 mal 12 Fächern:
\par

\begin{LetterGrid}
C & O & S & T & A & E & I & U & C & Z & H & I \\ \hline
H & U & S & S & V & L & N & I & N & E & E & E \\ \hline
R & Y & U & S & G & T & E & U & H & R & B & C \\ \hline
M & C & E & N & S & L & S & C & E & N & I & T \\ \hline
I & L & Ä & I & G & A & I & S & U & L & S & D \\ \hline
L & G & E & N & D & L & D & S & B & I & N & W \\ \hline
R & H & I & S & G & R & A & I & R & A & G & C \\ \hline
C & E & I & N & I & D & E & E & H & I & I & E \\ \hline
R & N & Ä & W & S & R & I & T & W & L & I & N \\ \hline
H & L & R & I & H & H & E & E & B & Z & U & I \\ \hline
U & E & N & C & E & G & S & I & E & T & H & A \\ \hline
I & G & S & H & S & A & E & H & H & S & N & N \\ \hline
\end{LetterGrid}
Die 24 Buchstaben
\par

(\Sperr{CHRMILRCRHUI IECTDWCENIAN})
\par

der linken und rechten Seitencolumne geben anagrammatisch angeordnet
\par

\Sperr{CANT R CRE HUI DRM J W NIHIL CEC},
\par

d. h. Cant R(egiomontanus) cre(ator) hui(us) dr(a)m(atis), J(ohann) W(olfgang) (Goethe) nihil cec(init); Kant aus Königsberg ist Schöpfer dieses Dramas, Johann Wolfgang Goethe hat nichts davon gedichtet! — Will ich etwa die Geheimschriftschlüssel der Baconianer parodiren? Wer könnte das! Man bleibt viel zu vernünftig, ungezwungen, nüchtern, wenn man in diesem Fach dilettirt; der unfreiwillige Humor, die ganze Verschraubtheit und Verstiegenheit dieser Köpfe lässt sich nicht nachahmen. Wenn man jenen Goethe-Kant-Scherz gegen die berühmte Titelblatt-Chiffre hält, worin sich Bacon als den „Graveur“ der Shakespeareschen Dramen zu erkennen giebt, so ist man in Gefahr, den Ernst beim Parodisten und die Parodie beim Ernstmeinenden zu suchen. Der Baconianer findet eine Widmung, ein Sonett in auffallend schlechtem Englisch abgefasst: ein Wink! und nun holt er mit allem Hocuspocus eines unbekehrbaren Déchiffreurs den eigentlichen Satz heraus, der noch viel gequälteres, unmöglicheres Englisch und im Übrigen so klar wie das Hexeneinmaleins ist. — Man thut Unrecht, die Unvernunft summarisch als Gegensatz des Vernünftigen zu verstehen; es giebt Gradunterschiede, und die Shakespeare-Baconfrage ist, wie der Spiritismus, vielleicht ein Gipfel im Hochgebirge der Absurdität. Man entschliesst sich, Jemandem eine, zwei, drei Voraussetzungen wider den \Foreign{bon sens} zuzugeben, man zwingt sein Gehirn, mit den Schwingungen eines Querkopfgehirns parallel zu schwingen — aber das hilft noch nichts; innerhalb des Unsinnigen kommt erst das Unsinnigste, die Narrheit potenzirt sich. Es giebt Systeme, die auf einem falschen Grunde wenigstens richtig aufgebaut sind, man tritt um den Preis eines \Foreign{sacrifizio dell’ intelletto} ein, darf dann aber wieder seine Vernunft gebrauchen. Und es giebt andere, die auf Schritt und Tritt den Unsinn nicht loswerden, Systeme, die es mit einer einmaligen Gehirnverrenkung keineswegs bewenden lassen; die Vernunft wird von unten auf gerädert — gliedweise und fingerweise verstümmelt!
\par

